% Faithful transcription — original French. Transcribed from the original first printing:
% Émile Picard, "Sur les intégrales de différentielles totales et sur une classe de surfaces
% algébriques", Note présentée par M. Hermite, Comptes rendus hebdomadaires des séances de
% l'Académie des sciences, tome 99, 2e semestre, séance du lundi 22 décembre 1884
% (Paris: Gauthier-Villars, 1884), pp. 1147–1149.
% Scan: Bibliothèque nationale de France / Gallica (ark:/12148/bpt6k3055h, vues f1147–f1149).
%
% \origpage uses the PRINTED page numbers. The note begins partway down p. 1147, below Henri
% Poincaré's note on total differentials; that preceding note is not transcribed, nor are the running
% heads, volume/issue notices ("C. R., 1884, 2e Semestre..."), feuilleton marks ("153"), or isolated
% page numbers. The subsequent note by M. Amigues beginning on p. 1149 is likewise omitted.
%
% Page boundary joins:
% - p. 1147 ends mid-sentence ("... présenter quelque"); p. 1148 opens immediately with no blank
%   line after \origpage{1148} ("intérêt; parmi ces applications...").
% - p. 1148 ends mid-sentence ("... est de degré $m$ et de"); p. 1149 opens immediately with no blank
%   line after \origpage{1149} ("genre $\frac{m + 2}{2}$, ...").
%
% QUOTATION MARKS. Following Comptes rendus convention for communications quoted across paragraphs,
% a single « opens the first paragraph on p. 1147, » heads each following paragraph on pp. 1148 and
% 1149, and » closes the communication at the end of p. 1149. Guillemets are set tight against the
% text (R18, R22).
%
% Cross-page decisions are recorded in notation.md alongside this file.
\documentclass{article}
\usepackage{readmasters}
\begin{document}

\origpage{1147}
\section*{Analyse mathématique. --- Sur les intégrales de différentielles totales et sur une classe de surfaces algébriques. Note de M.~E.~Picard, présentée par M.~Hermite.}

«Dans une Communication récente (\emph{Comptes rendus}, séance du 1\textsuperscript{er} décembre 1884), j'ai énoncé une proposition fondamentale relative aux intégrales de différentielles totales algébriques restant toujours finies, intégrales que je désigne pour cette raison sous le nom d'\emph{intégrales de première espèce}. Les applications de cette théorie, que je développe en ce moment dans un Mémoire étendu, pourront peut-être présenter quelque
\origpage{1148}intérêt; parmi ces applications, une des plus simples est relative aux surfaces algébriques dont les coordonnées peuvent s'exprimer par des fonctions uniformes quadruplement périodiques de deux paramètres. Je demande la permission d'énoncer ici une proposition générale concernant ces surfaces.

»Étant donnée une surface d'ordre $m$
\[
f(x, y, z) = 0,
\]
n'ayant que les singularités ordinaires considérées dans la Note citée, cherchons comment on pourra reconnaître si l'on peut exprimer les coordonnées $x$, $y$, $z$ d'un point quelconque par des fonctions quadruplement périodiques de deux paramètres, et cela de telle manière qu'à un point quelconque de la surface ne corresponde qu'un seul système de valeurs des deux paramètres, abstraction faite des multiples des périodes. Les conditions \emph{nécessaires et suffisantes} pour qu'il en soit ainsi peuvent être formulées de la manière suivante~:

»\emph{La surface proposée a une courbe double d'ordre $\frac{m(m - 4)}{2}$ et elle possède deux intégrales de différentielles totales de première espèce}
\[
\int \frac{B\,dx - A\,dy}{f'_{z}} \quad\text{et}\quad \int \frac{B_{1}\,dx - A_{1}\,dy}{f'_{z}},
\]
\emph{pour lesquelles le déterminant $AB_{1} - A_{1}B$ n'est pas identiquement nul.}

»Le nombre $m$ est nécessairement pair, et il résulte des conditions précédentes que la surface n'a pas de points doubles \emph{isolés}; de plus, le \emph{genre} de la surface est égal à l'unité.

»Les deux intégrales précédentes auront \emph{quatre} paires de périodes simultanées, et les deux équations aux différentielles totales
\[
\begin{aligned}
\frac{B\,dx - A\,dy}{f'_{z}} &= du, \\
\frac{B_{1}\,dx - A_{1}\,dy}{f'_{z}} &= dv
\end{aligned}
\]
donneront pour $x$ et $y$ des fonctions uniformes, quadruplement périodiques, de $u$ et $v$.

»Nous avons montré comment on pourrait reconnaître si une surface admet des intégrales de première espèce, et les trouver quand elles existent; le problème proposé est donc complètement résolu.

»Une section plane quelconque de la surface est de degré $m$ et de
\origpage{1149}genre $\frac{m + 2}{2}$, et cette courbe présente une particularité intéressante au point de vue de la réduction du nombre des périodes des intégrales abéliennes qui lui correspondent.

»Il résulte, en effet, immédiatement des résultats précédents que, parmi ces intégrales abéliennes, \emph{il y en a deux qui ont seulement quatre paires de périodes correspondantes}.»

\end{document}
